\section{Introduction}

In Monopoly, a deed that completes a rival's color group is worth little to its holder as property and a great deal as a veto. Whoever buys it gains a source of rent that every other player will pay, the seller included. Bargaining theory treats such a sale as one with identity-dependent externalities and predicts delay even when every party knows everything~\cite{jehiel1995delay}. When the holder later needs money, the buyer who values the deed most is the one rival it would complete, which is the standard setting for a forced sale below value~\cite{shleifer1992liquidation}. A four-player Monopoly game is therefore a small closed economy that puts to language-model agents the questions any multi-party market would. Do they notice what a trade gives a third party? Do they charge for it? What makes them give it up, and what do they do with the proceeds?

We examine these questions in twenty complete games among four configured LLM systems, abbreviated Sol, Gemini, Opus and Grok (Table~\ref{tab:config}). The agents buy, auction, trade, build, mortgage and go bankrupt through explicit legal actions, and they negotiate in free text through proposals and counteroffers. Each decision records the legal menu, the delivered state, every response attempt, a public message and a private note. The corpus holds \fact{track.prompts} decisions, \fact{barg.census.threads} negotiation threads, \fact{barg.census.responses-all.n} responses to offers, \fact{debt.episodes} debt episodes and \fact{debt.eliminated} eliminations. The study is observational. Its unit is the configured system in these games, including its response handling and the information it received.

Our central finding is that attention to rivals is common to all four systems, and liquidity is what separates their trajectories. Three results support it.

\emph{Every system treats a completing deed as the rival's claim on a monopoly.} Holders refuse most requests for such deeds or answer with a higher price. Asking prices rise as the buyer nears a complete group in every coarse cut we can form, although the \fact{d.pairs.n} cases in which one seller prices one deed for both kinds of buyer are mixed. Recipients also refuse trades that would complete their own group, usually with a stated consequence for the rival.

\emph{The deed moves when its holder is short of cash.} Holders with under \$100 reject \fact{d.gap.low.pct}\% of completion requests, and holders with \$300 or more reject \fact{d.gap.high.pct}\%. The gap appears in each of the five seed blocks. Where a sale followed an earlier refusal by the same holder, the holder's cash had fallen in between in \fact{barg.dyads.from-first-holder-decline.lower} of \fact{barg.dyads.from-first-holder-decline.n} cases.

\emph{What the cash becomes separates the trajectories.} Some cash-poor sellers were paying for their own houses, and most of them won. Others were postponing insolvency, and none of them won. After completing a group by trade, Sol and Gemini almost always built at once, while Opus and Grok often did not. Paying a debt that can be paid is the one financial task every system performed reliably.

The paper contributes a complete account of bargaining over rival-completing deeds in a multi-player LLM game, with every offer, the cash prices and \fact{barg.dyads.dyads} holder and buyer pairs followed to their resolution (Sections~\ref{sec:attention} and~\ref{sec:liquidity}). It describes, with an adjustment for exposure, what completed groups and sale proceeds became (Section~\ref{sec:conversion}). It reports the full population of debt episodes, paid and unpaid, and each game's position before its first estate changed hands (Section~\ref{sec:default}). Eight contextual incidents with source-exact quotations accompany the counts.

Four boundaries apply throughout. We score no decision against an optimal policy, because the cohort contains no value reference. A stated reason is a recorded explanation, and we do not read it as the cause of an action or as a measurement of internal computation~\cite{turpin2023unfaithful,lanham2023faithfulness}. The prompt invites bluffing and deception in public messages, so a public bluff is authorized conduct and we do not call it misalignment. Twenty games in five seed blocks with one cyclic seating order cannot rank four systems, and Opus played most of them under a different response condition (Section~\ref{sec:conditions}).

The project began from eight questions about LLM agents in strategic settings. They concern strategic reasoning, opponent modeling, negotiation and deception, long-horizon planning and memory, adaptation, economic decisions and risk, collective dynamics, and the validity of the evaluation itself. Three broader questions ask whether agents are strategic or only sound strategic (A), whether general capability carries over to multi-agent settings (B), and whether skill can be told apart from luck and weak opponents (C). Section~\ref{sec:discussion} states what this cohort says about each and what it leaves unmeasured.
